\begin{abstract}
\noindent
Мы строим и исследуем полный цикл генеративной рекомендательной системы на
датасете Yambda-50M с официальным временным разбиением: авторегрессионную
модель следующего айтема, reward model с четырьмя головами поведения и
офлайн-дообучение политики методом GRPO. В едином прогоне, в котором базовая
модель, reward model и RL обучены на одном коммите кода, базовая модель
(81.5 млн параметров, 4 слоя, контекст 200 событий, словарь 630\,974 айтема,
sampled softmax с 8192 негативами) достигает на тестовом дне NDCG@10 0.0282
и Recall@100 0.0659; по NDCG@10 она в 1.63 раза превосходит популярностный
baseline (0.0173). Reward model отличает лайк от случайного айтема каталога
(ROC-AUC 0.751), но на правдоподобных кандидатах из истории того же
пользователя сигнал почти исчезает (ROC-AUC 0.564), а долю досмотра она
ранжирует слабо (Spearman 0.07). Обучение с подкреплением за 64 шага
повышает награду модели на 4.4\% и снижает NDCG@10 на 82\%: парный бутстрэп
по пользователям даёт $\Delta = -0.0230$ и 95\% CI $[-0.0254, -0.0206]$
для NDCG@10 и $\Delta = -0.0471$, CI $[-0.0505, -0.0438]$ для Recall@100.
Мы разбираем механизм падения: топ-10 разных пользователей схлопывается
(Jaccard $0.030 \to 0.736$), средняя длина рекомендованного трека падает с
181 до 124 с, а оценка головы «лайк» в выдаче уменьшается вдвое при росте
оценки головы «досмотр». Это переоптимизация суррогатной награды с измеренным
механизмом: политика увеличивает ровно ту голову, которую оптимизирует, за
счёт остальных. Отдельный прогон с базовой моделью, обученной на 30\,000
шагов, даёт NDCG@10 0.0623; на нём эффект RL не измерялся. Работа фиксирует
направление эффекта на одном seed и определяет границы, в которых её выводы
имеют силу.
\end{abstract}
